\scanpage{201}

Notons que l'on peut aussi considérer

\[
f_{X/S}:X\mathbin{\unsure{-|\longrightarrow}}X^{p/S}
\]

d'où, en itérant, des

\[
f_{X/S}^{p^i/S}:X^{p^i/S}\longrightarrow X^{p^{i+1}/S}
\]

et par suite un composé

\[
F_{X/S}^{(\nu)}=f_{X/S}^{(p^{\nu-1})}f_{X/S}^{(p^{\nu-2})}\ldots f_{X/S}^{\deleted{\unsure{p}}}:
X\mathbin{-}X^{(p^\nu/S)}
\]

et on aura, si $X/S$ est \uline{lisse} \uline{de dim relative} $1$

\[
F_*^{(\nu)}=\operatorname{Tr}_{F_{X/S}^{(\nu)}}=\underset{\text{Gysin}}{\bigl(F_{X/S}^{\nu}\bigr)_*}
\]

en identifiant (sous réserve de formule de chgt. de base) $\deleted{\unsure{X^{p^i}}}$\deleted{\unsure{$f_*$}} $g_{\nu*}(\underline\Omega^1_{X^{p^\nu/S}})$ avec $g_{0*}(\Omega^1_{X/S})^{(p^{\unsure{\nu|1}})}=\underline\omega^{(p^\nu)}$. On trouve en particulier, si $\lambda\longmapsto\lambda^{p^\nu}$ est l'identité sur $\mathcal O_S^{*}$, donc $X^{p^\nu/S}\overset{\mathrm{can}}{\simeq}X$, que

\[
F^\nu_{X/S}:X\mathbin{-}X
\]

est un $S$-endomorphisme, qui d'ailleurs n'est autre que l'itéré $\nu$\textsuperscript{ième} de $F_X^\nu$ (compte tenu que par hyp. $F_S^\nu=\mathrm{id}_S$, \unsure{de sorte|du fait} que $F_X^\nu$ est bien un $S$-endomorphisme) [En général, et sans hyp. sur $S$, on peut dire que le composé

\diagram{\begin{tikzcd}[ampersand replacement=\&]
X\arrow[r,"F_{X/S}^{(\nu)}"] \& X^{(p^\nu/S)}\arrow[r,"\mathrm{can}"]\arrow[d,dotted] \& X\arrow[d,dashed]\\
\& {\bigl[\,S}\arrow[r,"{F_S^\nu}"] \arrow[ur,phantom,"\mathrm{cart}" description] \& {S\,\bigr]}
\end{tikzcd}}

est $F_X^\nu$ ]

% archive sheet number: 200
